\section{\texorpdfstring{Capítulo 1: Problema del Bandido de
\emph{k}-brazos}{Capítulo 1: Problema del Bandido de k-brazos}}\label{capuxedtulo-1-problema-del-bandido-de-k-brazos}

\subsection{1.1. Introducción}\label{introducciuxf3n}

\subsubsection{Descripción del Problema y
Relevancia}\label{descripciuxf3n-del-problema-y-relevancia}

El problema del bandido de \(k\)-brazos (o \emph{Multi-Armed Bandit})
ilustra formalmente el dilema fundamental del Aprendizaje por Refuerzo:
el balance iterativo entre la exploración y la explotación. Un agente
computacional se enfrenta a un entorno estacionario y estocástico
modelado como una máquina tragaperras con \(k\) palancas (brazos), cada
una ocultando una distribución de probabilidad de recompensa subyacente
diferente y desconocida \emph{a priori}.

Con un tiempo límite o presupuesto de \(T\) interacciones, la meta del
agente es maximizar su recompensa acumulativa. La toma de su decisión es
crítica: si confía únicamente en la palanca que le ha concedido mayor
ganancia económica hasta ahora (explotación), recae en el riesgo de
converger hacia una recompensa local ignorando una palanca numéricamente
superior. Si, por el contrario, invierte un esfuerzo desmedido en probar
nuevas palancas aleatorias para descubrir el verdadero margen de
beneficio (exploración), desperdicia iteraciones cobrando recompensas
bajas con la consecuente penalización temporal.

\subsubsection{Motivación del Trabajo}\label{motivaciuxf3n-del-trabajo}

La razón principal por la que estudiamos este problema es porque se
parece a muchas decisiones prácticas que tomamos en el mundo real (como
elegir en qué anuncio invertir dinero o qué tratamiento médico probar).
Es un entorno ``simple'' porque las decisiones que tomamos ahora no
cambian las opciones que tendremos en el futuro (cada tirada es
independiente). Sin embargo, estudiarlo es clave para entender cómo
funciona la inteligencia artificial cuando aprende a base de prueba y
error.

Si logramos comprender cómo estos algoritmos básicos
(\(\epsilon\)-Greedy o UCB) interactúan con la duda de \emph{``¿me quedo
con lo que ya sé que funciona o arriesgo a probar algo nuevo?''},
habremos construido los pilares necesarios para entender a los
verdaderos agentes robóticos complejos que resolverán problemas mucho
más difíciles en el futuro.

\subsubsection{Objetivos del Informe}\label{objetivos-del-informe}

El desarrollo de esta primera fase trata de cumplir los siguientes
objetivos técnicos: 1. Diseñar y programar un entorno modular en Python
capaz de modelar un Bandido bajo tres tipologías de distribución de
recompensa continuas y discretas: Normal, Binomial y Bernoulli. 2.
Implementar desde cero tres familias clásicas de algoritmos de toma de
decisiones: Explotación estocástica (\(\epsilon\)-Greedy y sus variantes
con decaimiento), Ascenso de Gradiente (Softmax y métricas de
Preferencia) y Optmismo basado en incertidumbre (UCB-1, UCB-2 y
UCB1-Tuned). 3. Evaluar dichas familias sometiéndolas empíricamente
frente al entorno mediante \(N=500\) ejecuciones independientes
simulando \(T=300\) pasos temporales. 4. Analizar la escalabilidad, la
ratio de acierto de selección óptima y la mitigación del arrepentimiento
acumulado (\emph{Cumulative Regret}).

\subsubsection{Organización del
Documento}\label{organizaciuxf3n-del-documento}

Para facilitar su asimilación, este primer capítulo estructura su
narrativa en el siguiente formato secuencial. En la sección de
\textbf{Desarrollo} (1.2), se sentarán las bases teóricas definiendo
matricial y probabilísticamente las distribuciones a las que se enfrenta
el bandido, adjuntando a la par un breve estado del arte de las
influencias bibliográficas. Seguidamente, en la sección de
\textbf{Algoritmos} (1.3), se expondrá la matemática formal y el
pseudocódigo subyacente de cada política heurística testeada. Las
asunciones y las simulaciones se materializan descriptivamente de la
mano de gráficas de rendimiento en la sección referente a la
\textbf{Evaluación/Experimentos} (1.4). Por último, las métricas son
discutidas críticamente en las \textbf{Conclusiones} (1.5) para
certificar el descubrimiento de los parámetros ganadores globales.

\subsection{1.2. Desarrollo}\label{desarrollo}

\subsubsection{Contexto y Antecedentes}\label{contexto-y-antecedentes}

El dilema de exploración frente a explotación es una encrucijadas
elemental dentro del Aprendizaje por Refuerzo. Mientras que en el
aprendizaje supervisado un algoritmo recibe explícitamente la respuesta
correcta a modo de etiqueta, en el aprendizaje evaluativo
(característico del RL) un agente debe descubrir autónomamente qué
acciones brindan la mayor ganancia probándolas. El modelo del
\emph{Multi-Armed Bandit} (Bandido Multibrazo) es el paradigma para
aislar y estudiar estas propiedades evaluativas y de decisión en un
entorno estacionario y no asociativo (es decir, donde el entorno no
cambia de estado). Comprender la balanza entre arriesgar ganancias
seguras presentes y el potencial de ganancias futuras desconocidas
cimienta la matemática que los agentes complejos aplican cuando operan
con los valores \(Q\) y variables de estado multidimensionales.

\subsubsection{Definición Formal del
Problema}\label{definiciuxf3n-formal-del-problema}

El entorno estocástico se define bajo el espacio de recompensas
discretas o continuas para un conjunto finito de acciones o brazos
\(a \in \mathcal{A} = \{1, 2, \dots, k\}\).

El verdadero y desconocido ``valor esperado'' de elegir una acción \(a\)
se denota habitualmente como el valor a maximizar \(q_*(a)\), y está
formulado como la esperanza matemática de la recompensa en un instante
\(t\): \[ q_*(a) = \mathbb{E}[R_t \mid A_t = a] \]

A lo largo del presente marco de validación, las recompensas \(R_t\)
derivan su valor estocástico basándose en tres entornos implementados:

\begin{enumerate}
\def\labelenumi{\arabic{enumi}.}
\item
  \textbf{Entorno de Recompensas Normales (Gausianas)} La retribución
  estocástica de cada acción se muestrea de una distribución diferencial
  \(R_t \sim \mathcal{N}(\mu_a, \sigma_a^2)\), donde la verdadera medida
  subyacente de recompensa está dominada por su función de densidad de
  probabilidad (PDF):
  \[ f(x \mid \mu_a, \sigma^2) = \frac{1}{\sqrt{2\pi\sigma^2}} e^{-\frac{(x - \mu_a)^2}{2\sigma^2}} \]
\item
  \textbf{Entorno de Recompensas Binomiales} En este escenario, la
  activación del brazo \(a\) devuelve una recompensa entera discreta
  (\(x \in \mathbb{N}\)) que contabiliza la cantidad total de éxitos
  sobre \(n_a\) ensayos independientes. Es dirigido por la variable de
  masa probabilística \(R_t \sim B(n_a, p_a)\):
  \[ P(R_t=x) = \binom{n_a}{x} p_a^x (1 - p_a)^{n_a - x} \]
\item
  \textbf{Entorno de Recompensas de Bernoulli} Un escenario más sencillo
  donde el resultado de tirar de la palanca solo puede ser un éxito o un
  fracaso absoluto (\(R_t \in \{0, 1\}\)). Modela la variable
  \(R_t \sim Bernoulli(p_a)\) en donde el valor promedio de iteración a
  estabilizar es exactamente la probabilidad intrínseca al brazo
  evaluado (\(q_*(a) = p_a\)):
  \[ P(R_t) = p_a^{R_t} (1 - p_a)^{1 - R_t} \]
\end{enumerate}

\subsubsection{Enfoque y Justificación
Metodológica}\label{enfoque-y-justificaciuxf3n-metodoluxf3gica}

Para abordar el problema estocástico, se ha abstraído un
\emph{framework} orientado a objetos que desacopla completamente el
``Entorno'' del ``Agente''. Esto permite insertar módulos de políticas
que basan su decisión algorítmica interactuando con las mismas
estructuras pre-configuradas \(R_t\). El presente trabajo se caracteriza
por implementar la comparativa simultánea y simétrica del estado del
arte (\emph{\(\epsilon\)-Greedy} frente al Límite de Confianza
\emph{UCB} y el cálculo heurístico \emph{Softmax}) bajo un volumen
estadístico (\(N=500\) iteraciones promediadas), garantizando un estudio
numérico de optimización con mitigación del \emph{Arrepentimiento
Acumulado} en ecosistemas tanto lógicos continuos como discretos.

\subsubsection{1.2.1. Trabajos
Relacionados}\label{trabajos-relacionados}

El planteamiento algorítmico, su formulación paramétrica y las
heurísticas estandarizadas analizadas en este proyecto se fundamentan y
se derivan de la literatura existente en el campo del Aprendizaje por
Refuerzo: - \textbf{Métodos Epsilon y Ascenso de Gradiente (Sutton y
Barto, 2018):} El diseño de las estrategias \(\epsilon\)-Greedy y de los
métodos probabilísticos como el \emph{Ascenso de Gradiente} se ha
construido siguiendo las reglas y fundamentos explicados en el libro de
referencia ``Reinforcement Learning: An Introduction''
\cite{sutton2018reinforcement}. - \textbf{Manejo de la Incertidumbre con
UCB (Auer et al., 2002):} Para programar los algoritmos UCB (Upper
Confidence Bound) de forma robusta y entender cómo utilizan sus fórmulas
matemáticas para explorar las opciones menos conocidas, nos hemos
apoyado en el estudio ``Finite-time Analysis of the Multiarmed Bandit
Problem'' \cite{auer2002finitetime}.

\subsection{1.3. Algoritmos}\label{algoritmos}

Para resolver el problema del bandido de \(k\)-brazos hemos implementado
tres enfoques distintos: métodos que combinan explotar lo conocido con
explorar al azar, algoritmos que premian probar opciones inciertas para
dejar de dudar, y tácticas que eligen las acciones basándose en
preferencias relativas en lugar de en ganancias directas. A continuación
se justifican sus enfoques y se detalla el funcionamiento lógico de los
pilares principales de cada familia algorítmica estudiada.

\subsubsection{\texorpdfstring{Familia 1: Lógicas de
Exploración-Explotación
(\(\epsilon\)-Greedy)}{Familia 1: Lógicas de Exploración-Explotación (\textbackslash epsilon-Greedy)}}\label{familia-1-luxf3gicas-de-exploraciuxf3n-explotaciuxf3n-epsilon-greedy}

\textbf{Justificación:} El algoritmo \(\epsilon\)-Greedy es el punto de
partida más clásico en el Aprendizaje por Refuerzo. Su funcionamiento es
muy sencillo y rápido de calcular: la mayor parte del tiempo (con
probabilidad \(1-\epsilon\)) el agente elige la opción que hasta ahora
le ha dado más beneficios, pero de vez en cuando (con una pequeña
probabilidad \(\epsilon\)) elige una opción al azar para explorar. A
partir de aquí nace una mejora llamada \(\epsilon\)-Decay, que soluciona
el problema de explorar para siempre: va reduciendo poco a poco la
probabilidad de explorar (multiplicándola por un factor \(\lambda\)), de
modo que, a medida que el agente aprende, se vuelve más seguro y se
centra en asegurar las ganancias.

\textbf{Pseudocódigo de \(\epsilon\)-Greedy Clásico:}

\begin{algorithm}
\caption{Algoritmo $\epsilon$-Greedy para K-Brazos}
\begin{algorithmic}[1]
\REQUIRE $\epsilon \in [0, 1]$, Total de acciones $K$
\STATE Inicializar $Q(a) \leftarrow 0$ para todo $a=1,\dots, K$
\STATE Inicializar $N(a) \leftarrow 0$ para todo $a=1,\dots, K$
\FOR{cada paso $t=1, 2, \dots, T$}
    \STATE $p \leftarrow$ Valor aleatorio de Distribución Uniforme(0, 1)
    \IF{$p < \epsilon$}
        \STATE $A_t \leftarrow$ Acción aleatoria uniforme en $\{1, \dots, K\}$
    \ELSE
        \STATE $A_t \leftarrow \arg\max_a Q(a)$  \COMMENT{Ruptura de empates al azar}
    \ENDIF
    \STATE Ejecutar $A_t$, observar recomensa $R_t$
    \STATE $N(A_t) \leftarrow N(A_t) + 1$
    \STATE $Q(A_t) \leftarrow Q(A_t) + \frac{1}{N(A_t)} \big[R_t - Q(A_t)\big]$
\ENDFOR
\end{algorithmic}
\end{algorithm}

\subsubsection{Familia 2: Optimismo ante la Incertidumbre
(UCB)}\label{familia-2-optimismo-ante-la-incertidumbre-ucb}

\textbf{Justificación:} A diferencia de explorar siempre al azar, los
métodos de \emph{Upper Confidence Bound} (UCB) son curiosos de forma
mucho más inteligente. UCB calcula cuánta confianza tiene en lo que ya
sabe de cada palanca. Si una opción se ha probado muy poco, su nivel de
incertidumbre es alto, así que el algoritmo decide elegirla para ``salir
de dudas''. Esto convierte a UCB en un método robusto y bueno para no
perder el tiempo. En este trabajo hemos estudiado el algoritmo básico
UCB1 (que usa un número \(c\) para controlar de forma fija esa
curiosidad), el modelo UCB2 (que optimiza el proceso probando las
opciones en ``ráfagas'' o bloques cada vez más largos en lugar de
recalcular a cada paso), y la variante avanzada UCB1-Tuned (que se
ajusta a sí misma midiendo la variabilidad de los premios que recibe).

\textbf{Pseudocódigo de UCB2:}

\begin{algorithm}
\caption{Algoritmo UCB2}
\begin{algorithmic}[1]
\REQUIRE Constante exploratoria $\alpha \in (0, 1)$
\STATE Inicializar contador de épocas por brazo: para cada $a$, $r_a \leftarrow 0$
\STATE Iniciar explorando todos los brazos una vez, actualizando su recompensa promedio $Q(a)$
\FOR{paso de tiempo $t = K+1, \dots$}
    \STATE Seleccionar brazo $A_t \leftarrow \arg\max_a \left[ Q(a) + \sqrt{\frac{(1 + \alpha) \ln (e \cdot t / \tau(r_a))}{2 \tau(r_a)}} \right]$ \newline
           (donde $\tau(r) = \lceil (1+\alpha)^r \rceil$)
    \STATE Determinar longitud de la ráfaga: $\Delta \leftarrow \tau(r_{A_t} + 1) - \tau(r_{A_t})$
    \FOR{iteraciones $i = 1$ hasta $\Delta$}
        \STATE Ejecutar $A_t$, percibir recompensa $R$
        \STATE $N(A_t) \leftarrow N(A_t) + 1$
        \STATE $Q(A_t) \leftarrow Q(A_t) + \frac{1}{N(A_t)} \big[R - Q(A_t)\big]$
    \ENDFOR
    \STATE Actualizar época del brazo: $r_{A_t} \leftarrow r_{A_t} + 1$
\ENDFOR
\end{algorithmic}
\end{algorithm}

\subsubsection{Familia 3: Ascenso de Gradiente Estocástico (Softmax y
Preferencias)}\label{familia-3-ascenso-de-gradiente-estocuxe1stico-softmax-y-preferencias}

\textbf{Justificación:} Si los algoritmos solo se fijan en los puntos
exactos de cada opción (\(Q\)), pueden tener problemas para decidir si
los premios son muy parecidos (por ejemplo, ganar \(+85\) o \(+84\)).
Los métodos basados en el Gradiente lo resuelven usando ``Preferencias''
(\(H_a\)). En vez de memorizar valores numéricos rígidos, simplemente
aprenden qué opción prefieren en relación al resto y deciden usando
porcentajes (con la Ley Exponencial Suave o Gibbs). En nuestras pruebas,
esta estrategia del \emph{Gradiente de Preferencias} trataremos de
demostrar si es la que mejor se adapta a los entornos donde las
recompensas son simplemente ``todo o nada'' (como sucede en el modelo de
\emph{Bernoulli}).

\textbf{Pseudocódigo de Base Gibbs / Gradiente:}

\begin{algorithm}
\caption{Algoritmo de Variación de Gradiente de Preferencias}
\begin{algorithmic}[1]
\REQUIRE Tasa de aprendizaje $\alpha > 0$
\STATE Inicializar preferencias $H(a) \leftarrow 0 \;\; \forall a$ 
\STATE Inicializar recompensa promedio de validación $\bar{R} \leftarrow 0$
\FOR{$t = 1, \dots, T$}
    \STATE Calcular probabilidades de la curva Softmax:\\
    $P(a) \leftarrow \frac{\exp(H(a))}{\sum_{b=1}^K \exp(H(b))} \quad \forall a$
    \STATE Elegir acción $A_t$ en base a la distribución de densidad $P(\cdot)$
    \STATE Obtener recompensa $R_t$
    \STATE Actualizar la base empírica global $\bar{R} \leftarrow \bar{R} + \frac{1}{t}(R_t - \bar{R})$
    \FORALL{$a \in \{1, \dots, K\}$}
        \IF{$a = A_t$}
            \STATE $H(a) \leftarrow H(a) + \alpha(R_t - \bar{R})(1 - P(a))$
        \ELSE
            \STATE $H(a) \leftarrow H(a) - \alpha(R_t - \bar{R})P(a)$
        \ENDIF
    \ENDFOR
\ENDFOR
\end{algorithmic}
\end{algorithm}

\subsection{1.4. Evaluación y
Experimentos}\label{evaluaciuxf3n-y-experimentos}

\subsubsection{1.4.1. Configuración
Experimental}\label{configuraciuxf3n-experimental}

Para garantizar la validez científica y la total reproducibilidad de los
resultados a lo largo del estudio de las tres familias de distribuciones
(Normal, Binomial y Bernoulli), el entorno de pruebas fue estandarizado
con los siguientes hiperparámetros de simulación global: -
\textbf{Ejecuciones (\emph{Runs}):} \(N = 500\) pruebas independientes
estocásticas. Las curvas finales presentadas corresponden a la matriz
promediada probabilísticamente sobre este número total de semillas. -
\textbf{Horizonte de Simulador (\emph{Steps}):} \(T = 300\) secuencias
temporales de decisión por ejecución. - \textbf{Semilla Base Estática:}
Se inyectó globalmente una semilla matemática (aleatoria constante,
habitualmente \texttt{np.random.seed(1024)}) para inicializar la
topología generatriz del bandido de forma idéntica ante la comparativa
de los tres algoritmos.

Para llevar a cabo las simulaciones, se definieron los siguientes
espectros de hiperparámetros iniciales para cada modelo inteligente
(Tabla \ref{tab:Parametros}):

\begin{table*}[t]
\centering
\caption{Hiperparámetros definidos y comparados por familia algorítmica.}
\label{tab:Parametros}
\begin{tabular}{|l|l|l|}
\hline
\textbf{Familia Algorítmica} & \textbf{Configuración / Algoritmo} & \textbf{Hiperparámetros Evaluados} \\ \hline
Exploración-Explotación & $\epsilon$-Greedy & $\epsilon \in \{0.0, 0.01, 0.1\}$ \\ \hline
Exploración-Explotación & $\epsilon$-Decay & Tasa de decaimiento $\lambda \in \{0.8, 0.99, 0.9999\}$ \\ \hline
Ascenso de Gradiente & Softmax & Temperatura $\tau \in \{0.5, 1.0, 2.0\}$ \\ \hline
Ascenso de Gradiente & Gradiente de Preferencias & Tasa de aprendizaje $\alpha \in \{0.5, 1.0, 2.0\}$ \\ \hline
Límite Sup. Condianza & UCB1 & Constante de exploración $c \in \{0.5, 1.0, 1.5\}$ \\ \hline
Límite Sup. Condianza & UCB2 & Expansión de confianza temporal $\alpha \in \{0.25, 0.5, 0.75\}$ \\ \hline
Límite Sup. Condianza & UCB1-Tuned & Ninguno (Varianza de la Recompensa empírica) \\ \hline
\end{tabular}
\end{table*}

\subsubsection{1.4.2. Métricas de
Desempeño}\label{muxe9tricas-de-desempeuxf1o}

La evaluación del rendimiento integral de los agentes propuestos se
analizó multidimensionalmente observando su desempeño iterativo bajo 4
prismas métricos: 1. \textbf{Recompensa Promedio Continua (\emph{Average
Reward}):} Traza iterativa de la estimación \(\bar{R_t}\); refleja la
capacidad del agente de percibir y adaptar ganancia bruta promedio. 2.
\textbf{Ratio Probabilístico de Selección (\emph{\% Optimal Action}):}
Discrimina de forma binaria el acierto o el fallo explícito al localizar
matemáticamente la mejor máquina subyacente \(\arg\max q_*(a)\). 3.
\textbf{Arrepentimiento Acumulado (\emph{Cumulative Regret}):} Mide la
acumulación de pérdida teórica u ``oportunidad perdida'' a nivel
temporal de decantarse por opciones penalizadoras frente a un escenario
``clarividente''. Se busca priorizar gráficas cóncavas y rápidamente
estabilizadas (planas). 4. \textbf{Histogramas Distribucionales:}
Estadísticos puros del muestreo frecuentista con el comportamiento
latente final del ecosistema (\emph{Arm Statistics}).

\subsubsection{1.4.3. Análisis de Resultados (Gráficas y
Comparativas)}\label{anuxe1lisis-de-resultados-gruxe1ficas-y-comparativas}

La convergencia de rendimiento obtenida tras los experimentos numéricos
extraídos de \texttt{epsilon\_greedy.ipynb},
\texttt{ascenso\_gradiente.ipynb} y \texttt{UCB\_algorithms.ipynb} ha
posibilitado la categorización de parámetros óptimos para las tres ramas
del estado del arte, englobados en la siguiente matriz global (Tabla
\ref{tab:Resultados}):

\begin{table*}[t]
\centering
\caption{Parámetros con mayor rendimiento cruzados por distribución.}
\label{tab:Resultados}
\begin{tabular}{|l|l|p{7cm}|}
\hline
\textbf{Familia} & \textbf{Mejor Hiperparámetro} & \textbf{Análisis Crítico} \\ \hline
$\epsilon$-Greedy   & $\epsilon = 0.1$                                  & Globalmente superior a versiones estáticas conservadoras ($\epsilon=0.01$). \\ \hline
Decaimento      & $\lambda = 0.99$                                  & Combina convergencia exponencial de exploración. Cae veloz frente a entornos como Bernoulli si su suelo ($\epsilon_{min}$) es severo.\\ \hline
Softmax         & Temperatura $\tau \in [0.5, 1.0]$                 & Curva modesta en continuas; catastrófico (20\% acierto) en dominios Bernoulli. \\ \hline
P. Gradiente    & Ratio de Aprendizaje $\alpha \in [0.5, 1.0]$      & La revelación. Su matriz de distancias absolutas mitiga por completo entornos acotados y discretos (Bernoulli).  \\ \hline
UCB             & $c = 1.0$ (o $c=0.5$ Bernoulli)                    & Curvas de convergencia teórica muy seguras y logarítmicas. Rinde maravillosamente para incertidumbres Gausianas. \\ \hline
UCB1-Tuned      & Ninguno (No paramétrico)                          & Defensa inquebrantable en todas las distribuciones; el entorno escala su nivel de exploración por la varianza natamente. \\ \hline
\end{tabular}
\end{table*}

\subsubsection{\texorpdfstring{Resultados Familia
\(\epsilon\)-Greedy}{Resultados Familia \textbackslash epsilon-Greedy}}\label{resultados-familia-epsilon-greedy}

\begin{itemize}

\item
  \textbf{Decaimiento (\(\epsilon\)-Decay):} Aunque en teoría debería
  ser mejor, en la práctica el decaimiento no ha supuesto ninguna
  ventaja clara en los entornos \emph{Normal} y \emph{Binomial}. De
  hecho, ha funcionado igual o ligeramente \textbf{peor} que el modelo
  básico (\(\epsilon\)-Greedy puro). De todos los ajustes probados,
  reducir la probabilidad de explorar progresivamente con un
  multiplicador \(\lambda=0.99\) fue lo más equilibrado. En cambio, si
  la exploración se reduce muy despacio (\(\lambda=0.9999\)), el
  algoritmo pierde demasiado tiempo escogiendo opciones al azar,
  acumulando muchas pérdidas.
\item
  \textbf{Bernoulli}: En este caso particular, hemos experimentado un
  problema grave: si frenamos la caída de exploración demasiado pronto
  (fijando un mínimo inamovible de \(\epsilon_{min}=0.001\)), el
  algoritmo deja de aprender de golpe. Esto provoca que se estanque
  prematuramente y no logre igualar la enorme tasa de aciertos que sí
  consiguen otros competidores.
\end{itemize}

\subsubsection{Resultados Familia de Ascenso de
Gradiente}\label{resultados-familia-de-ascenso-de-gradiente}

\begin{itemize}

\item
  \textbf{Softmax frente a los entornos:} Mientras que el método Softmax
  se defiende de forma aceptable en entornos donde los premios varían de
  forma continua (distribución \emph{Normal}), fracasa completamente
  cuando los premios son de ``todo o nada'' (\emph{Bernoulli}). En este
  último caso, el algoritmo se atasca rápidamente, eligiendo la mejor
  opción apenas un 20\% de las veces y acumulando muchísimas pérdidas a
  lo largo del tiempo.
\item
  \textbf{Gradiente de Preferencias vs Softmax:} Este ha sido uno de los
  grandes descubrimientos de las pruebas. El \emph{Gradiente de
  Preferencias}, al guiarse por ``cuánto le gusta'' una opción
  (porcentajes de preferencia) en lugar de intentar memorizar valores
  exactos de puntos, \textbf{ha superado con mucha claridad a Softmax en
  todos los escenarios}. Su éxito es especialmente espectacular en el
  entorno \emph{Bernoulli}, donde consigue elegir la mejor palanca casi
  el 100\% de las veces desde muy temprano. Esto hace que sus pérdidas
  (arrepentimiento) se queden en unos ridículos 3 puntos tras 300
  intentos (frente a los 70 puntos perdidos por Softmax).
\end{itemize}

\subsubsection{Resultados Familia UCB}\label{resultados-familia-ucb}

\begin{itemize}

\item
  \textbf{UCB1 vs.~UCB2:} El modelo clásico UCB1 aprende bastante bien
  en casi todas las situaciones probadas, sobre todo cuando su nivel de
  confianza inicial está fijado en \(c=1.0\) (para los modelos
  \emph{Normal} y \emph{Binomial}). Sin embargo, en el juego de premios
  binarios (\emph{Bernoulli}) le va mejor si arranca más precavido con
  un \(c=0.5\). Su hermano mayor, \textbf{UCB2}, muestra dos caras
  opuestas: \textbf{falla estrepitosamente en los entornos de Binomial y
  Bernoulli}, donde se atasca eligiendo constantemente la peor opción
  sin aprender casi nada, pero a cambio resulta ser rapidísimo y el
  mejor de todos cuando los premios son más naturales y variados
  (entornos \emph{Normales}).
\item
  \textbf{UCB1-Tuned}: Esta versión con ``piloto automático'' consiguió
  igualar algunos de los buenos resultados de UCB1 sin necesidad de
  configurar ningún parámetro a mano. Aún así, sigue pasándolo bastante
  mal cuando las reglas del juego son drásticas (como en el entorno
  \emph{Bernoulli}), donde necesita atascarse al menos 100 intentos
  seguidos antes de lograr darse cuenta de cuál es el brazo ganador.
\end{itemize}

Para ver todo esto de forma mucho más clara, las siguientes gráficas
(sacadas de probar los algoritmos en entornos de premios de tipo
\emph{Normal}) muestran un resumen visual perfecto. Tal y como
descubrimos al analizar a fondo las simulaciones, \textbf{la gráfica del
Arrepentimiento Acumulado resulta ser siempre la más útil y fácil de
entender}. Esto se debe a que, a diferencia de otras medidas, en el
Arrepentimiento se puede distinguir a simple vista si un algoritmo sigue
aprendiendo y reduciendo sus pérdidas o si ya se ha estancado por
completo.

\begin{figure}[h]
    \centering
    \includegraphics[width=\columnwidth]{k_brazos/img/eg_regret_normal.png}
    \caption{Curva de Arrepentimiento Acumulado para la familia $\epsilon$-Greedy.}
    \label{fig:egreedy_perf}
\end{figure}

En la Figura \ref{fig:egreedy_perf}, se discierne claramente cómo una
caída suave del decaimiento temporal (\(\lambda = 0.99\)) permite
achatar asintóticamente el \emph{regret} mitigando estancamientos
lineales que sí sufrían ramas no decadentes.

\begin{figure}[h]
    \centering
    \includegraphics[width=\columnwidth]{k_brazos/img/ag_regret_normal.png}
    \caption{Curva de Arrepentimiento Acumulado comparando Preferencias vs Softmax.}
    \label{fig:ag_perf}
\end{figure}

Por otro lado, la comparación que vemos en la Figura \ref{fig:ag_perf}
nos demuestra algo muy interesante: aunque el algoritmo Softmax arranca
bien y consigue ganancias rápidamente, basar las decisiones puramente en
las \textbf{Preferencias} de cada opción frente al resto (la curva
verde) resulta ser una estrategia absolutamente imbatible a medio y
largo plazo siempre que se tenga configurada la tasa de aprendizaje
correcta (\(\alpha=1\)).

\begin{figure}[h]
    \centering
    \includegraphics[width=\columnwidth]{k_brazos/img/ucb_regret_normal.png}
    \caption{Curva de Arrepentimiento Acumulado demostrando el inquebrantable éxito de UCB1.}
    \label{fig:ucb_perf}
\end{figure}

Por último, la Figura \ref{fig:ucb_perf} demuestra en la práctica lo
bien que funciona la idea principal del método UCB: si castigamos
fuertemente a las opciones que ya hemos probado demasiadas veces,
conseguimos que el algoritmo frene en seco sus pérdidas
(arrepentimiento). Todo esto lo logra con una precisión asombrosa, de
forma rápida y sin necesidad de fórmulas ni ajustes complicados, tal y
como atestigua el excelente recorrido de la variante con ``piloto
automático'' (\emph{UCB1-Tuned} en la curva verde).

\subsection{1.5. Conclusiones}\label{conclusiones}

\subsubsection{Resumen de Principales
Resultados}\label{resumen-de-principales-resultados}

Tras probar a fondo todos los algoritmos en el problema del bandido de
k-brazos, sacamos una conclusión clara: no existe un algoritmo perfecto
para todo. La mejor estrategia siempre dependerá del tipo de entorno y
de cómo se comporten los premios.

\begin{itemize}

\item
  En entornos donde las recompensas varían de forma continua y lógica
  (como en la distribución \emph{Normal}), métodos como \textbf{UCB1} y
  especialmente \textbf{UCB2} aprenden rapidísimo al elegir de forma
  inteligente qué opciones explorar. Sin embargo, \textbf{solo UCB1}
  logra mantener ese buen nivel en entornos algo más rígidos como el
  \emph{Binomial}, ya que UCB2 se atasca por completo.
\item
  Cuando nos enfrentamos a escenarios drásticos de ``todo o nada''
  (recompensas \emph{Bernoulli}), casi todos los algoritmos anteriores
  fallan de forma estrepitosa y acumulan pérdidas. En estos casos tan
  extremos, el ganador es el \textbf{Gradiente de Preferencias}. En vez
  de obsesionarse con calcular exactamente cuántos puntos da cada
  opción, este modelo se limita a comparar qué tan buena es una palanca
  respecto de las demás; gracias a esto, logra dominar el juego y no
  equivocarse casi nunca (\(Regret \approx 0\)).
\item
  Para situaciones variadas o cuando no tenemos mucha potencia de
  cálculo, recurrir a los clásicos métodos semialeatorios (como
  \textbf{\(\epsilon\)-Decay} o directamente un \(\epsilon\) fijo de
  \(0.1\)) resulta ser una estrategia buena.
\end{itemize}

\subsubsection{Limitaciones del Estudio}\label{limitaciones-del-estudio}

La limitación principal de los resultados reportados en este documento
es que el juego asume que las reglas no cambian, la propiedad de
\textbf{estacionalidad} del modelo K-Brazos abstraído. En el mundo real,
la característica estocástica y las medias de recompensa (\(\mu\)) de
una máquina o variable evolucionan y mutan con el paso del tiempo. Si
transladamos estos algoritmos a la vida real, estos algoritmos podrían
fracasar estrepitosamente porque se obcecarían en descartar opciones de
escenarios que de repente podrían volverse muy rentables.

Además, este marco experimentacional excluye variables de ``contexto''
(variables climáticas, visuales o previas a la decisión) de cada tirada
o instante \(t\), tratando a las palancas bajo el paradigma
probabilístico absoluto e idéntico en cada \emph{step}.

\subsubsection{Impacto y Reflexión en el Aprendizaje por
Refuerzo}\label{impacto-y-reflexiuxf3n-en-el-aprendizaje-por-refuerzo}

Estudiar a fondo estas tres familias de algoritmos del bandido (apostar
a lo seguro, investigar las opciones menos conocidas, o guiarse por las
preferencias) no es para nada un ejercicio anticuado o puramente
teórico. Al contrario: el concepto del \emph{Arrepentimiento} y cómo
minimizarlo son los cimientos más básicos que le dicen a cualquier
Inteligencia Artificial cómo debe tomar sus decisiones.

Estas fórmulas e ideas matemáticas conforman el motor interno inicial de
decisiones de cualquier agente moderno y complejo en el mundo real (como
los sistemas de \emph{Deep Reinforcement Learning} de las IA generativas
o de los coches autónomos).

\subsubsection{Líneas Futuras de
Estudio}\label{luxedneas-futuras-de-estudio}

El siguiente paso natural para avanzar en este campo sería crear
simulaciones de \textbf{Bandidos No Estacionarios (\emph{Non-stationary
Bandits})}, es decir, enseñar a la Inteligencia Artificial a jugar en
entornos donde las reglas y los premios van cambiando con el tiempo.
Para lograrlo, habría que dotar a los algoritmos con capacidades de
adaptación ``olvidadizas'': mecanismos matemáticos que les permitan
ignorar lecciones antiguas que ya no sirven y centrarse en lo que está
ocurriendo ahora mismo.

Más adelante, el objetivo final sería evolucionar hacia \textbf{Bandidos
Contextuales (\emph{Contextual Bandits})}. En este escenario, la IA no
jugaría a ciegas, sino que recibiría ``pistas'' del entorno antes de
tener que elegir una opción (por ejemplo, saber qué tiempo hace antes de
recomendar un tipo de ropa). Esto abre directamente las puertas hacia
sistemas de inteligencia artificial mucho más completos y complejos que
los que hemos visto hasta ahora.
